\section{Introduction}
\label{sec:introduction}

Text-SPARQL example sets---paired collections of natural language questions and executable SPARQL queries---are central to modern Knowledge Graph Question Answering (KGQA). They serve as in-context demonstrations that steer large language model (LLM) query generation~\cite{wikisp2023,zahera2024generating,emonet2024llm,dabramo-etal-2025-investigating}, as training corpora for fine-tuned translation models~\cite{sgpt2022generative,neural2022machine}, and as standardized benchmarks for evaluating system accuracy~\cite{trivedi2017lc,dubey2019lcquad,usbeck2019qald,usbeck2021qald}. However, high-quality datasets remain scarce: existing benchmarks are concentrated almost exclusively on a handful of prominent open-domain knowledge graphs (principally DBpedia, Wikidata, and Freebase) where substantial human curation and community effort have been invested~\cite{usbeck2019qald,dubey2019lcquad,liu2024spinach}. In contrast, domain-specific knowledge graphs in specialized disciplines such as biomedicine~\cite{kgs2025biomedical} and geoscience~\cite{sockg2025} lack such community infrastructure and possess virtually no Text-SPARQL resources, severely impeding the development, adaptation, and evaluation of KGQA systems in real-world domains.

Automatically synthesizing Text-SPARQL datasets for arbitrary knowledge graphs without human authoring presents three fundamental technical challenges that naive LLM prompting cannot overcome:
\begin{enumerate}
    \item \textbf{Schema Hallucination under Cold-Start Access:} Real-world knowledge graphs often feature extensive, evolving, or poorly documented schemas. When presented only with a SPARQL endpoint, LLMs frequently hallucinate non-existent classes and properties.
    \item \textbf{Multi-Hop Relational Validity (Empty Executions):} LLMs cannot infer multi-hop graph connectivity in the dark. Unguided query synthesis routinely generates structurally disconnected triple patterns that return empty result sets ($\emptyset$) upon execution.
    \item \textbf{Syntactic and Semantic Homogeneity:} Naive prompting collapses into repetitive, single-hop entity-attribute lookups, failing to capture the structural diversity (e.g., value filtering, aggregations, optional clauses, property paths) required for robust benchmarking.
\end{enumerate}

To overcome these challenges, we present \textbf{QueryBloom}, an automated, KG-agnostic framework that synthesizes diverse, execution-verified Text-SPARQL example sets from any RDF knowledge graph given only SPARQL endpoint access. QueryBloom is grounded in the observation that every SPARQL query represents an executable specification of an underlying subgraph pattern. Rather than generating queries top-down from unconstrained prompts, QueryBloom inverts this relationship: it first discovers structurally sound, multi-hop subgraph patterns directly from the knowledge graph and grounds them with concrete entity instances and attribute values. These grounded subgraphs then seed a multi-agent distillation pipeline that generates natural language questions, injects syntactic complexity (such as filters, unions, and optional patterns), and iteratively verifies execution results against the live endpoint to guarantee structural validity, non-empty execution, and semantic alignment.

We evaluate QueryBloom by synthesizing datasets over DBpedia and comparing them against the human-authored DBpedia QALD-9 benchmark across three key dimensions: query structural complexity, semantic linguistic variety, and holistic schema coverage. Our empirical findings show that QueryBloom-generated datasets match or exceed established human-authored benchmarks: QueryBloom queries exhibit four times the rate of \texttt{FILTER} constraints as QALD-9 and successfully incorporate advanced constructs such as \texttt{OPTIONAL} blocks and property paths that are absent from the baseline. Furthermore, sentence embedding analysis demonstrates broader semantic question coverage, while schema metrics reveal three times deeper triple patterns and a substantially higher ratio of unique query templates.

In summary, our key contributions are:
\begin{itemize}
    \item \textbf{KG-Agnostic Synthesis Methodology:} We introduce QueryBloom, an automated framework that combines endpoint-driven multi-hop subgraph sampling with multi-agent LLM distillation and execution-guided verification to synthesize high-quality Text-SPARQL pairs given only SPARQL endpoint access.
    \item \textbf{Addressing Domain Scarcity:} We provide a general solution to the cold-start data scarcity bottleneck in domain-specific knowledge graphs, enabling automated benchmark and training data creation for arbitrary RDF graphs without manual ontology engineering.
    \item \textbf{Comprehensive Empirical Validation:} We conduct an in-depth evaluation demonstrating that QueryBloom-synthesized datasets surpass human-crafted benchmarks in structural query complexity, semantic question diversity, and schema coverage.
\end{itemize}
